% Archived earlier experiment-design draft.
% The current thesis structure uses chapters/04_experimental_setup_and_evaluation.tex.
\chapter{Experimental Design and Evaluation}
\label{ch:archived_experiment_design}

The Cartesian impedance-control strategy developed in this thesis is
evaluated with respect to three criteria:

\begin{enumerate}
  \item \textbf{Stability} --- Does the closed-loop system remain stable
        for all tested stiffness and damping parameters?
        Can stability be guaranteed analytically?
  \item \textbf{Tracking behavior} --- How accurately and quickly does the
        end-effector converge to the desired position and orientation after a step change?
        What are the settling time and steady-state error as functions of
        the stiffness matrix $\mathbf{K}_p$?
  \item \textbf{Compliant interaction} --- Does the controller correctly
        implement the virtual wall constraint?
        Is the end-effector free to slide within the plane and to tilt
        until the tool face makes full contact with the surface,
        while being stiff against normal penetration and spinning?
\end{enumerate}

The theoretical basis for all three criteria is developed in
\cref{ch:theory}.
This chapter translates those criteria into four concrete experiments:

\begin{itemize}
  \item \textbf{Experiment~1} (\cref{sec:exp_tracking}): fixed-point tracking with rotated stiffness directions.
  \item \textbf{Experiment~2} (\cref{sec:exp_wall}): Virtual wall
        constraint --- normal stiffness, in-plane sliding, tilt-to-contact.
  \item \textbf{Experiment~3} (\cref{sec:exp_stiffness}): Normal stiffness
        variation --- validates the compliance law $d_{ss}=F_n/k$.
  \item \textbf{Experiment~4} (\cref{sec:exp_config}): Configuration
        dependence --- same Cartesian forces and moments at three different arm configurations.
\end{itemize}

% ─────────────────────────────────────────────────────────────────
\section{Preliminary: Stability Analysis}
\label{sec:exp_stability}
% ─────────────────────────────────────────────────────────────────

\subsection{Lyapunov Stability Proof}

The closed-loop stability of the Cartesian impedance controller is
established using a \textbf{Lyapunov argument} \cite{Siciliano2009}.
Consider the total mechanical energy of the virtual mass-spring-damper
system as a candidate Lyapunov function:

\begin{equation}
\label{eq:lyapunov}
V(\mathbf{e}_p, \dot{\mathbf{p}}_{\mathrm{EE}}) =
\underbrace{\tfrac{1}{2}\,\mathbf{e}_p^\top\,\mathbf{K}_p\,\mathbf{e}_p}_{\text{potential energy}}
+
\underbrace{\tfrac{1}{2}\,\dot{\mathbf{p}}_{\mathrm{EE}}^\top\,\mathbf{M}_p\,\dot{\mathbf{p}}_{\mathrm{EE}}}_{\text{kinetic energy}}
> 0
\end{equation}

$V > 0$ for all $(\mathbf{e}_p, \dot{\mathbf{p}}_{\mathrm{EE}}) \neq (\mathbf{0},\mathbf{0})$
because both $\mathbf{K}_p$ and $\mathbf{M}_p$ are \abbr{SPD}.
This is precisely why \abbr{SPD} is required for both matrices: positive definiteness
of $\mathbf{K}_p$ ensures the potential energy is zero only at the desired
position, and positive definiteness of $\mathbf{M}_p$ ensures the kinetic
energy is zero only when the end-effector is at rest.

Taking the time derivative along the system trajectories and substituting the
positional impedance law \cref{eq:positional_impedance_force}:

\begin{align}
\dot{V}
&= \mathbf{e}_p^\top\,\mathbf{K}_p\,\dot{\mathbf{e}}_p
 + \dot{\mathbf{p}}_{\mathrm{EE}}^\top\,\mathbf{M}_p\,\ddot{\mathbf{p}}_{\mathrm{EE}}
\notag\\
&= -\,\mathbf{e}_p^\top\,\mathbf{K}_p\,\dot{\mathbf{p}}_{\mathrm{EE}}
 + \dot{\mathbf{p}}_{\mathrm{EE}}^\top\bigl(\mathbf{K}_p\,\mathbf{e}_p
   - \mathbf{D}_p\,\dot{\mathbf{p}}_{\mathrm{EE}}\bigr)
\notag\\
\label{eq:Vdot}
&= -\,\dot{\mathbf{p}}_{\mathrm{EE}}^\top\,\mathbf{D}_p\,\dot{\mathbf{p}}_{\mathrm{EE}}
\;\leq\; 0
\end{align}

where $\dot{\mathbf{e}}_p = -\dot{\mathbf{p}}_{\mathrm{EE}}$ (error decreases when
the end-effector moves towards $\mathbf{p}_d$).
The spring terms cancel exactly, leaving only the damping dissipation.

\paragraph{Why $\mathbf{D}_p$ must be SPD.}
$\dot{V} \leq 0$ requires the quadratic form
$\dot{\mathbf{p}}_{\mathrm{EE}}^\top\mathbf{D}_p\dot{\mathbf{p}}_{\mathrm{EE}} \geq 0$.
Strict inequality ($\dot{V} < 0$ for $\dot{\mathbf{p}}_{\mathrm{EE}} \neq \mathbf{0}$)
requires $\mathbf{D}_p$ to be positive definite.
This is the formal reason why a damping matrix with negative or zero
eigenvalues would destabilise the controller: it would allow $\dot{V}>0$
along some velocity direction, meaning the system gains energy in that
direction.

\paragraph{Conclusion.}
Since $V > 0$ and $\dot{V} \leq 0$, Lyapunov's theorem guarantees that
the equilibrium $(\mathbf{e}_p, \dot{\mathbf{p}}_{\mathrm{EE}}) = (\mathbf{0},\mathbf{0})$ is
\textbf{stable}.
Asymptotic stability (convergence to zero) follows from LaSalle's
invariance principle: the only trajectory along which $\dot{V} = 0$
permanently (i.e.\ $\dot{\mathbf{p}}_{\mathrm{EE}} \equiv \mathbf{0}$) is the
equilibrium itself.
The same argument applies to the rotational subsystem with $\mathbf{e}_R$,
$\mathbf{K}_R$, $\boldsymbol{\omega}_{\mathrm{EE}}$, $\mathbf{D}_R$.

\begin{remark}
This proof holds for \emph{any} SPD $\mathbf{K}_p$ and $\mathbf{D}_p$,
including the full off-diagonal matrices arising from the virtual wall
task.
The diagonal special case is not required for stability.
\end{remark}

\subsection{Experimental Stability Verification}

To confirm stability experimentally, the following protocol is executed:

\begin{enumerate}
  \item The robot is held at the desired position $\mathbf{p}_d$.
  \item The controller is activated with stiffness values ranging from
        $k_{\min}=50\,\mathrm{N/m}$ to $k_{\max}=1000\,\mathrm{N/m}$.
  \item A step disturbance is applied: the robot is manually displaced
        by $\Delta p = 30\,\mathrm{mm}$ and released.
  \item The system must return to equilibrium without sustained oscillation
        for the selected damping values.
  \item Joint torques and end-effector position are recorded.
        Divergence, sustained oscillation, or saturation of joint torques
        at their limits would indicate instability.
\end{enumerate}

\textbf{Success criterion:} The position error $\|\mathbf{e}_p(t)\|$
converges monotonically to zero for all tested values of $\mathbf{K}_p$
with the selected damping values.


For later comparison, the critically damped single-axis reference response is described by

\begin{equation}
\label{eq:settling}
t_s \approx \frac{4}{\omega_n}=4\sqrt{\frac{M_i}{K_i}},
\qquad \omega_n=\sqrt{\frac{K_i}{M_i}},
\end{equation}

where \(t_s\) is the approximate 2\,\% settling time and \(K_i\) is the
stiffness in the investigated principal direction for the corresponding
directional mass \(M_i=\Lambda_{ii}\). Under
a constant external force \(F_i\), the steady-state error is

\begin{equation}
\label{eq:ss_error}
e_{p,\infty}=\frac{F_i}{K_i}.
\end{equation}

These two relations are used to interpret the tracking and stiffness-variation
experiments.

\cref{fig:lyapunov} shows the simulated Lyapunov function $V(t)$ and
its time derivative $\dot{V}(t)$ for a 50\,mm step error with
$K=200\,\mathrm{N/m}$.
$V(t)$ decreases monotonically and $\dot{V}(t)\leq0$ throughout,
confirming the stability proof analytically.

\begin{figure}[htb]
\centering
\includegraphics[width=\textwidth]{lyapunov.pdf}
\caption{Stability verification: Lyapunov function $V(t)$ (left) and its
         time derivative $\dot{V}(t)=-\mathbf{D}_p\dot{p}_{\mathrm{EE}}^2\leq0$ (right)
         for a 50\,mm step error with $K=200\,\mathrm{N/m}$.
         The monotone decrease of $V$ and the non-positive $\dot{V}$ confirm
         asymptotic stability.}
\label{fig:lyapunov}
\end{figure}

% ─────────────────────────────────────────────────────────────────
\section{Experiment~1 --- Fixed-Point Tracking with Rotated Stiffness Directions}
\label{sec:exp_tracking}
% ─────────────────────────────────────────────────────────────────

This experiment evaluates the tracking behaviour of the impedance controller
for a single fixed Cartesian target point. In contrast to a purely diagonal
axis-aligned test, the stiffness and damping matrices are intentionally
constructed in a rotated local stiffness frame and then transformed into the
base frame using \cref{eq:positional_stiffness_base,eq:positional_damping_base}. This experiment
therefore demonstrates that full matrices are meaningful even when the task is
only to track one fixed point.

\subsection{Objective}

The objective is to show that the desired point \(\mathbf{p}_d\) defines the
centre of the virtual spring, while \(\mathbf{K}_{p,\mathrm{base}}\) defines
the shape and orientation of the restoring behaviour around that point. If the
principal stiffness directions are rotated relative to the base frame and the
stiffness values are not equal, the base-frame stiffness matrix becomes full.

\subsection{Task Setup}

The desired fixed point is chosen as

\begin{equation}
\label{eq:exp1_pd}
\mathbf{p}_d =
\begin{bmatrix}
0.50\\0.00\\0.40
\end{bmatrix}
\mathrm{m}.
\end{equation}

The initial position is offset by \(10\,\mathrm{mm}\) in the base \(x\)-direction:

\begin{equation}
\label{eq:exp1_initial}
\mathbf{p}_{EE}(0) =
\begin{bmatrix}
0.49\\0.00\\0.40
\end{bmatrix}
\mathrm{m}.
\end{equation}

Thus, the initial position error is

\begin{equation}
\label{eq:exp1_error}
\mathbf{e}_p(0)
=
\mathbf{p}_d-\mathbf{p}_{\mathrm{EE}}(0)
=
\begin{bmatrix}
0.01\\0\\0
\end{bmatrix}
\mathrm{m}.
\end{equation}

\subsection{Rotated Stiffness Field}

The local stiffness frame is rotated relative to the base frame using

\begin{equation}
\label{eq:exp1_W}
\mathbf{W}_{\mathrm{local}\rightarrow\mathrm{base}}
=
\mathbf{R}_z(40^\circ)\mathbf{R}_y(30^\circ)\mathbf{R}_x(20^\circ).
\end{equation}

The local stiffness and damping matrices are

\begin{equation}
\label{eq:exp1_KD_local}
\mathbf{K}_{p,\mathrm{local}}
=
\operatorname{diag}(1000,300,100)\;\mathrm{N/m},
\qquad
\mathbf{D}_{p,\mathrm{local}}
=
2\,\operatorname{diag}(\sqrt{1000},\sqrt{300},\sqrt{100}).
\end{equation}

They are transformed into the base frame as

\begin{equation}
\label{eq:exp1_KD_base}
\mathbf{K}_{p,\mathrm{base}}
=
\mathbf{W}_{\mathrm{local}\rightarrow\mathrm{base}}
\mathbf{K}_{p,\mathrm{local}}
\mathbf{W}_{\mathrm{local}\rightarrow\mathrm{base}}^{T},
\qquad
\mathbf{D}_{p,\mathrm{base}}
=
\mathbf{W}_{\mathrm{local}\rightarrow\mathrm{base}}
\mathbf{D}_{p,\mathrm{local}}
\mathbf{W}_{\mathrm{local}\rightarrow\mathrm{base}}^{T}.
\end{equation}

The expected base-frame stiffness matrix is approximately

\begin{equation}
\label{eq:exp1_Kbase_expected}
\mathbf{K}_{p,\mathrm{base}}
\approx
\begin{bmatrix}
605.1 & 359.0 & -205.6\\
359.0 & 500.2 & 138.9\\
-205.6 & 138.9 & 294.7
\end{bmatrix}
\;\mathrm{N/m}.
\end{equation}

\subsection{Expected Behaviour}

At the beginning of the experiment the robot is momentarily at rest, so the
initial damping force is zero. The initial restoring force is therefore

\begin{equation}
\label{eq:exp1_initial_force}
\mathbf{f}(0)
=
\mathbf{K}_{p,\mathrm{base}}\mathbf{e}_p(0)
\approx
\begin{bmatrix}
6.05\\3.59\\-2.06
\end{bmatrix}
\mathrm{N}.
\end{equation}

This is the key result of the experiment. Although the initial error is only in
the base \(x\)-direction, the restoring force contains \(x\)-, \(y\)-, and
\(z\)-components. This demonstrates that the full stiffness matrix defines a
rotated spring field around the fixed target point.

\subsection{Protocol}

\begin{enumerate}
  \item Move the end-effector to the initial position given in
        \cref{eq:exp1_initial}.
  \item Define the desired point \(\mathbf{p}_d\) according to
        \cref{eq:exp1_pd}.
  \item Construct \(\mathbf{K}_{p,\mathrm{base}}\) and
        \(\mathbf{D}_{p,\mathrm{base}}\) using
        \cref{eq:exp1_KD_base}.
  \item Activate the impedance controller and record
        \(\mathbf{p}_{\mathrm{EE}}(t)\), \(\dot{\mathbf{p}}_{\mathrm{EE}}(t)\),
        \(\mathbf{f}(t)\), and \(\boldsymbol{\tau}(t)\) at 1\,kHz.
  \item Compare the measured initial force direction with the predicted
        vector in \cref{eq:exp1_initial_force}.
\end{enumerate}

\subsection{Measured Variables and Expected Plots}

The following quantities are evaluated:
\begin{itemize}
  \item position error components \(e_{p,x}(t),e_{p,y}(t),e_{p,z}(t)\);
  \item Cartesian force components \(f_x(t),f_y(t),f_z(t)\);
  \item joint torques \(\tau_1(t),\ldots,\tau_7(t)\);
  \item comparison between diagonal stiffness and rotated full stiffness.
\end{itemize}

A diagonal axis-aligned stiffness matrix would generate only an \(x\)-force for
the initial condition in \cref{eq:exp1_error}. The rotated full matrix generates
a coupled 3D restoring force. This difference is the central observation of the
experiment.


% ─────────────────────────────────────────────────────────────────
\section{Experiment~2 --- Virtual Wall Constraint}
\label{sec:exp_wall}
% ─────────────────────────────────────────────────────────────────

This experiment evaluates whether the Cartesian impedance controller can create
an anisotropic interaction behaviour. The end-effector should remain close to a
virtual wall, resist motion in the wall-normal direction, and remain compliant
in tangential directions. The task is formulated using the local/base matrix
notation introduced in \cref{subsec:KD_matrices}.

\subsection{Protocol}

\begin{enumerate}
  \item The robot is moved to the initial position and orientation on the wall surface
        ($d = 0$, desired orientation $\mathbf{R}_d$ such that the
        end-effector $z$-axis aligns with $\hat{\mathbf{n}}$).
  \item The impedance controller is activated.
  \item A human operator applies forces by hand in three separate trials:
    \begin{itemize}
      \item \textbf{Trial A --- Normal push:}
            Force applied along $\hat{\mathbf{n}}$.
            Expected result: small displacement ($<1\,\mathrm{mm}$),
            large restoring force.
      \item \textbf{Trial B --- In-plane slide:}
            Force applied along $\hat{\mathbf{t}}_1$ or $\hat{\mathbf{t}}_2$.
            Expected result: end-effector slides freely in the plane;
            no normal displacement.
      \item \textbf{Trial C --- Tilt to contact:}
            Tool starts at an angle to the surface (touching at one point).
            The controller is activated with free tilt
            ($k_{R,\mathrm{free}}$ about $\hat{\mathbf{t}}_1$,
            $\hat{\mathbf{t}}_2$).
            Expected result: the tool tilts freely until its face is
            parallel to the surface (full contact); spinning about
            $\hat{\mathbf{n}}$ is resisted.
    \end{itemize}
  \item Joint torques, end-effector position, and orientation are recorded
        at 1\,kHz throughout each trial.
\end{enumerate}

\subsection{Simulated Results}

The three trials are simulated using the single-degree-of-freedom
impedance model with the parameter values from the previous subsection.
\Cref{fig:exp2_wall_trials} shows the displacement and rotation angle
for each trial.

\begin{figure}[htb]
\centering
\includegraphics[width=\textwidth]{exp2_wall_trials.pdf}
\caption{Simulation of the three virtual wall trials.
         \textbf{Top (Trial A):} a 10\,N normal push causes a
         $d_{\max}\approx9.8\,\mathrm{mm}$ displacement; after release
         the end-effector returns to $d=0$ within $0.13\,\mathrm{s}$.
         \textbf{Middle (Trial B):} a 2\,N in-plane force causes free
         sliding ($k_{\mathrm{free}}=50\,\mathrm{N/m}$) while the
         normal displacement remains near zero.
         \textbf{Bottom (Trial C):} a 0.5\,Nm roll torque rotates
         the tool freely ($k_{R,\mathrm{free}}=5\,\mathrm{Nm/rad}$)
         while spinning about $\hat{\mathbf{n}}$ is blocked ($k_{R,\mathrm{stiff}}=200\,\mathrm{Nm/rad}$).}
\label{fig:exp2_wall_trials}
\end{figure}

\begin{figure}[htb]
\centering
\includegraphics[width=\textwidth]{exp2_joint_torques.pdf}
\caption{Joint torques during Trial A (10\,N normal push at $t=0.5$--$1.5\,\mathrm{s}$).
         Each torque $\tau_k = J_{k,z}\cdot F_z$ is proportional to the
         $z$-component of the $k$-th Jacobian column at the home
         configuration.
         Joint~3 contributes most strongly; joint~7 contributes zero
         because its rotation axis is perpendicular to the $z$-force.}
\label{fig:exp2_joint_torques}
\end{figure}

\subsection{Measured Variables}

\begin{itemize}
  \item \textbf{Normal displacement
        \(d(t)=\hat{\mathbf{n}}^\top(\mathbf{p}_{EE}-\mathbf{p}_0)\)}:
        should remain near zero during in-plane trials and return to zero
        after a normal push.
  \item \textbf{In-plane position} $\mathbf{p}_{EE} \cdot \hat{\mathbf{t}}_1$
        and $\mathbf{p}_{EE} \cdot \hat{\mathbf{t}}_2$ vs.\ time:
        should show free motion during in-plane trials.
  \item \textbf{Tilt angles} about $\hat{\mathbf{t}}_1$ and $\hat{\mathbf{t}}_2$
        extracted from $R_{\mathrm{EE}}$ vs.\ time: should converge freely to
        the angle that brings the tool face parallel to the surface.
  \item \textbf{Spin angle} about $\hat{\mathbf{n}}$ vs.\ time: should
        remain near zero ($< 2^{\circ}$) in all trials.
  \item \textbf{Joint torques} $\tau_1, \ldots, \tau_7$ vs.\ time:
        should reveal which joints contribute most to normal stiffness
        vs.\ in-plane compliance.
\end{itemize}

% ─────────────────────────────────────────────────────────────────

\section{Experiment~3 --- Stiffness Variation in the Normal Direction}
\label{sec:exp_stiffness}
% ─────────────────────────────────────────────────────────────────

The normal stiffness $k_{\mathrm{stiff}}$ is varied while keeping the
in-plane and rotational parameters fixed.
A constant force of approximately $5\,\mathrm{N}$ is applied along
$\hat{\mathbf{n}}$ by the operator.
The steady-state displacement $d_{\mathrm{ss}}$ is measured for three values:

\begin{center}
\renewcommand{\arraystretch}{1.35}
\begin{tabular}{@{}lll@{}}
\toprule
$k_{\mathrm{stiff}}$ [N/m] & Expected $d_{\mathrm{ss}}$ [mm]
  & Behaviour \\
\midrule
$200$  & $25.0$ & Compliant: wall yields under load \\
$500$  & $10.0$ & Moderate stiffness \\
$1000$ & $5.0$  & Stiff: wall barely moves \\
\bottomrule
\end{tabular}
\end{center}

This experiment validates the linear compliance law
$d_{\mathrm{ss}} = F_n / k_{\mathrm{stiff}}$ predicted by the impedance model
(\cref{eq:ss_error}).
\Cref{fig:exp2_stiffness,fig:compliance_law} show the simulated
results and the theoretical compliance curve.

\begin{figure}[htb]
\centering
\includegraphics[width=\textwidth]{exp2_stiffness_variation.pdf}
\caption{Stiffness variation experiment.
         \textbf{Left:} Time response of the normal displacement for three
         stiffness values ($F_n=5\,\mathrm{N}$ applied at $t=0.5\,\mathrm{s}$).
         Higher stiffness gives faster response and smaller steady-state
         displacement.
         \textbf{Right:} Comparison of theoretical ($F/k$) and simulated
         steady-state displacements --- agreement is within 1\,\%.}
\label{fig:exp2_stiffness}
\end{figure}

\begin{figure}[htb]
\centering
\includegraphics[width=0.72\textwidth]{compliance_law.pdf}
\caption{Compliance law $d_{ss}=F_n/k_{\mathrm{stiff}}$ for four force
         levels.
         Red dots mark the three stiffness values tested experimentally.
         The linear relationship between compliance and inverse stiffness
         is the defining characteristic of the impedance model.}
\label{fig:compliance_law}
\end{figure}

% ─────────────────────────────────────────────────────────────────
\section{Experiment~4 --- Configuration Dependence of Joint Torques}
\label{sec:exp_config}
% ─────────────────────────────────────────────────────────────────

The virtual wall task is repeated at three different base configurations
$\vq^{(1)}, \vq^{(2)}, \vq^{(3)}$ while keeping the same Cartesian
end-effector position and orientation and the same stiffness parameters.
For a fixed Cartesian force $\mathbf{F}$, the joint torques
$\vtau = \mathbf{J}^\top(\vq)\mathbf{F}$ differ because the Jacobian
$\mathbf{J}$ depends on the configuration.

The torque distribution across the seven joints is recorded for each
configuration and compared as a bar chart, illustrating how redundancy
allows the same Cartesian task to be distributed differently across the
joints depending on arm configuration.

% ─────────────────────────────────────────────────────────────────
\section{Sensitivity Analysis}
\label{sec:sensitivity}

A sensitivity analysis investigates how variations in the controller
parameters affect the system response, providing insight into the
robustness of the design and the consequences of parameter
mis-specification.

\subsection{Sensitivity to Stiffness}

The steady-state error under a constant force $F_n$ along the normal
direction is $e_{p,\infty} = F_n / k_{\mathrm{stiff}}$ (from
\cref{eq:ss_error}).
The sensitivity of the steady-state error to a relative change in
$k_{\mathrm{stiff}}$ is:

\begin{equation}
\frac{\partial e_{p,\infty}}{\partial k_{\mathrm{stiff}}}
= -\frac{F_n}{k_{\mathrm{stiff}}^2}
\end{equation}

At $k_{\mathrm{stiff}}=1000\,\mathrm{N/m}$ and $F_n=5\,\mathrm{N}$,
a $\pm10\%$ variation in stiffness changes the steady-state error by
$\pm0.05\,\mathrm{mm}$, which is negligible for the surface-following
task.
At $k_{\mathrm{stiff}}=200\,\mathrm{N/m}$, the same $\pm10\%$
variation changes the error by $\pm1.25\,\mathrm{mm}$, which may
be perceptible.
This shows that \textbf{low stiffness values are more sensitive to
parameter uncertainty} than high stiffness values.

\subsection{Sensitivity to Damping}

Under-damping relative to a selected stiffness can cause oscillatory
transients, whereas excessive damping can increase settling time without
improving steady-state accuracy. The critical-damping value
$D_i = 2\sqrt{M_iK_i}$ is therefore used only as an analytical
reference, where \(M_i=\Lambda_{ii}\) comes from the virtual Cartesian mass
matrix computed from the libfranka model. In the implemented controller, the damping values are selected
directly as controller parameters and the logged response is used to verify
that oscillations remain sufficiently small for the tested motions.

\subsection{Sensitivity to Model Errors}

The Coriolis compensation term $\bm{\tau}_c$ is computed from the libfranka
model, while gravity compensation for the external torque command is handled by
the Franka Control Interface.
Model errors in these compensation effects appear as a constant force offset in
the Cartesian space:

\begin{equation}
\delta\mathbf{F}_{\mathrm{model}} =
(\mathbf{J}^\top)^{-1}\,\delta\bm{\tau}_{\mathrm{model}}
\end{equation}

For the Panda without payload, manufacturer specifications state
a gravity model accuracy of $\pm0.5\,\mathrm{Nm}$ per joint.
Propagated through the Jacobian at the home configuration, this
corresponds to a worst-case Cartesian force offset of approximately
$0.3\,\mathrm{N}$, producing a position offset of
$0.3/k_{\mathrm{stiff}} = 0.3\,\mathrm{mm}$ at
$k_{\mathrm{stiff}}=1000\,\mathrm{N/m}$.
This is within the positioning tolerance and confirms that the
libfranka model is sufficiently accurate for the tasks considered.

\section{Offline Evaluation of Logged Data}
\label{sec:offline_evaluation}

The CSV file is evaluated offline using Python.
First, the positional and rotational error norms are computed from the logged error components.
The maximum error is used to identify the largest deviation during the experiment, while the steady-state error is computed as the mean error over the final 20\% of the experiment.
Overshoot is evaluated for each individual error component after the first zero crossing.
In addition, the maximum commanded torque of each joint is computed to verify that the torque commands remain within acceptable limits.
The logged signals are then plotted over time to analyze convergence, damping behavior, overshoot, and torque smoothness.

\section{Result Interpretation Guide}
\label{sec:exp_results}
% ─────────────────────────────────────────────────────────────────

The result interpretation connects every observation directly to the theory
and to one of the three evaluation criteria.
Example formulations organised by criterion:

\paragraph{Stability.}
\begin{itemize}
  \item \textit{``The position error $\|\mathbf{e}_p(t)\|$ converged
        monotonically to zero in all trials without oscillation, confirming
        the selected damping values used in the controller.''}
  \item \textit{``No saturation of the joint torque limits was observed,
        indicating that the stiffness values lie within the dynamically
        feasible range of the Panda actuators.''}
\end{itemize}

\paragraph{Tracking behavior.}
\begin{itemize}
  \item \textit{``The measured settling time at $K_x=200\,\mathrm{N/m}$
        was $t_s = 0.29\,\mathrm{s}$, consistent with the theoretical
        prediction $4/\omega_n = 4\sqrt{M_i/K_x}$ from
        \cref{eq:settling}, using the directional mass from the virtual
        Cartesian mass matrix.''}
  \item \textit{``The steady-state error under a $3\,\mathrm{N}$ external
        force at $K_x = 500\,\mathrm{N/m}$ was $6.1\,\mathrm{mm}$,
        matching the predicted value $F/K = 3/500 = 6.0\,\mathrm{mm}$.''}
\end{itemize}

\paragraph{Compliant interaction.}
\begin{itemize}
  \item \textit{``The normal displacement remained below $1\,\mathrm{mm}$
        during in-plane sliding (Experiment~2, Trial B), confirming that
        $k_{\mathrm{stiff}} = 1000\,\mathrm{N/m}$ is sufficient to maintain
        wall contact within the positioning accuracy of the sensor.''}
  \item \textit{``The roll angle increased freely to $\phi_{\mathrm{roll}} = 25^{\circ}$
        without any resistance above the noise floor, consistent with the
        prescribed $k_{R,\mathrm{free}} = 5\,\mathrm{Nm/rad}$.''}
  \item \textit{``Pitch and yaw remained within $\pm 1.5^{\circ}$ throughout all
        trials, confirming that $k_{R,\mathrm{stiff}} = 200\,\mathrm{Nm/rad}$
        prevents any spinning of the tool about the wall normal.''}
  \item \textit{``The torque distribution changed significantly across the
        three configurations even for the same Cartesian forces and moments, demonstrating
        the configuration dependence of the Jacobian transpose mapping.''}
\end{itemize}
